\subsection{Duality of modules of finite length}

Let A be a Noetherian ring; denote by \(\Omega\) the set of its maximal ideals. Generalizing the definition given in the preceding number, we shall say that an A-module J is a Matlis A-module if it is injective, if its associated prime ideals are the maximal ideals of A, and if for every maximal ideal \(\mathfrak{m}\) of A the A/ \(\mathfrak{m}\)-vector space \(\mathrm{Hom}_{\mathrm{A}}(\mathrm{A}/\mathfrak{m},\mathrm{J})\) has dimension 1. For each \(\mathfrak{m} \in \Omega\) , choose an injective envelope \(\kappa(\mathfrak{m}) \to \mathrm{I}(\mathfrak{m})\) of the A-module \(\kappa(\mathfrak{m})\); the A-module \(\bigoplus_{\mathfrak{m} \in \Omega} \mathrm{I}(\mathfrak{m})\) is a Matlis module, and every Matlis A-module is isomorphic to it ( \(n^{\circ}\) 1, th. 1).

Recall (VIII, § 1, no. 5) that we denote by \(Z_0(A)\) the \(Z\)-module \(Z^{(\Omega)}\) and \(\varepsilon : Z_0(A) \to Z\) the linear form which maps each element of the basis \(\Omega\) to 1. If M is an A-module of finite length, then \(A_m\) -module \(M_m\) is of finite length for every \(m \in \Omega\), and zero except for finitely many ideals \(m \in \Omega\). We set

\[
z _ {0} (M) = \sum_ {\mathfrak {m} \in \Omega} \operatorname{long} _ {\mathrm{A} _ {\mathfrak {m}}} (\mathrm{M} _ {\mathfrak {m}}) [ \mathfrak {m} ] \quad \text { in } \quad Z _ {0} (\mathrm{A});
\]

We have \(\mathrm{long}_{\mathrm{A}}(\mathrm{M}) = \varepsilon(z_0(\mathrm{M}))\) (loc. cit., example 3). Conversely, an A-module N such that \(\mathrm{long}_{\mathrm{A_m}}(\mathrm{N_m})\) is finite for every \(\mathfrak{m} \in \Omega\), and zero outside a finite subset I of \(\Omega\) , is of finite length: indeed N is isomorphic to a submodule of \(\bigoplus_{\mathfrak{m} \in I} \mathrm{N}_{\mathfrak{m}}\) (II, § 3, no. 3, cor. 2 of th. 1), and we have \(\mathrm{long}_{\mathrm{A_m}}(\mathrm{N_m}) = \mathrm{long}_\mathrm{A}(\mathrm{N_m})\) since every \(\mathrm{A_m}\) -simple module is isomorphic to \(\kappa(\mathfrak{m})\) , hence simple as an A-module.

Let J be a Matlis A-module. For every A-module M, we shall denote by \(D_{A}(M)\) or simply \(D(M)\) the A-module \(\operatorname{Hom}_{\mathrm{A}}(\mathrm{M},\mathrm{J})\). Let \(\alpha_{M}\) be the homomorphism from M to \(D(D(M))\) defined by \(\alpha_{M}(m)(f)=f(m)\) for \(m\in M\) , \(f\in D(M)\) .

PROPOSITION 4.--- For the A-module M to be of finite length, it is necessary and sufficient that D(M) be of finite length. We then have \(z_{0}(\mathrm{D}(\mathrm{M})) = z_{0}(\mathrm{M})\), \(\operatorname{long}_{\mathrm{A}}(\mathrm{M}) = \operatorname{long}_{\mathrm{A}}\mathrm{D}(\mathrm{M})\) , \(\operatorname{Ann}_{\mathrm{A}}(\mathrm{M}) = \operatorname{Ann}_{\mathrm{A}}(\mathrm{D}(\mathrm{M}))\) , and the A-linear map \(\alpha_{M}\) is bijective.

For every \(\mathfrak{m} \in \Omega\), the \(A_{\mathfrak{m}}\) -module \(D(M)_{\mathfrak{m}}\) is identified with \(\operatorname{Hom}_{A_{\mathfrak{m}}} (M_{\mathfrak{m}}, J_{\mathfrak{m}})\) (II, § 2, no. 7, prop. 19); the first assertion of the proposition then follows from th. 2, c) and from the characterization of modules of finite length given above. Suppose henceforth that M is of finite length; we have \(\operatorname{long}_{A_{\mathfrak{m}}} (D(M)_{\mathfrak{m}}) = \operatorname{long}_{A_{\mathfrak{m}}} (M_{\mathfrak{m}})\) for every \(\mathfrak{m} \in \Omega\) (loc. cit.), whence \(z_0(D(M)) = z_0(M)\) and \(\operatorname{long}_A(D(M)) = \operatorname{long}_A(M)\). Moreover, the map \((\alpha_{M})_{\mathfrak{m}}: M_{\mathfrak{m}} \to D(D(M))_{\mathfrak{m}}\) is identified with the canonical homomorphism \(\alpha_{M_{\mathfrak{m}}}\) , which is bijective (loc. cit.); consequently \(\alpha_{M}\) is bijective.

For every \(a \in A\), we have \(D(a_{M}) = a_{D(M)}\) and consequently \(\text{Ann}_{A}(M) \subset \text{Ann}_{A}(D(M))\). Applying this to the A-module \(D(M)\) we deduce the opposite inclusion, whence the equality \(\text{Ann}_{A}(M) = \text{Ann}_{A}(D(M))\) .

Example.--- Let A be a principal ideal ring, K its field of fractions. The A-module K/A is a Matlis module: indeed the canonical map from K/A to \(\prod_{\mathfrak{m}\in\Omega}\mathrm{K}/\mathrm{A}_{\mathfrak{m}}\) induces an isomorphism from K/A to \(\bigoplus_{\mathfrak{m}\in\Omega}\mathrm{K}/\mathrm{A}_{\mathfrak{m}}\) (A, VII, p. 10, th. 2); the assertion then follows from no. 1, example 1. We have moreover already proved in A, VII, § 4, no. 9 that the map \(\alpha_{M}\) is bijective for every A-module M of finite length when the ring A is principal.

